\documentclass[11pt,a4paper]{article}

\usepackage[utf8]{inputenc}
\usepackage[T2A]{fontenc}
\usepackage[russian,english]{babel}
\usepackage{amsmath,amssymb,amsthm,mathtools}
\usepackage{geometry}
\usepackage{booktabs}
\usepackage{enumitem}
\usepackage{hyperref}

\geometry{margin=2.5cm}

\title{Многоклассовая модель оптимального управления банком\\с различными дюрациями кредитов}
\author{}
\date{}

\newcommand{\1}{\mathbf{1}}
\newcommand{\E}{\mathbb{E}}
\newcommand{\R}{\mathbb{R}}
\newcommand{\pos}[1]{\left[#1\right]^+}

\begin{document}
\selectlanguage{russian}
\maketitle

\section{Общая структура баланса}

Рассматривается банк, предлагающий $M$ классов кредитов с фиксированной процентной ставкой и различными средними сроками погашения (дюрациями).
Баланс банка в момент времени $t$ имеет вид
\begin{equation}
    \sum_{i=1}^{M} A_{i,t} \;+\; C_t \;=\; D_t \;+\; K_t,
\end{equation}
где
\begin{itemize}[noitemsep]
    \item $A_{i,t}$ --- объём кредитного портфеля $i$-го класса,
    \item $C_t$ --- денежные средства (cash, резервы),
    \item $D_t$ --- привлечённые депозиты (с плавающей ставкой),
    \item $K_t$ --- собственный капитал.
\end{itemize}

Как и в базовой одноклассовой модели, предполагается, что капитал полностью размещён в денежных средствах ($C_t = K_t$), а объём депозитов всегда равен совокупному объёму кредитного портфеля:
\begin{equation}
    D_t = \sum_{i=1}^{M} A_{i,t} \equiv A_t.
\end{equation}
Таким образом, капитал выполняет роль буфера ликвидности, не участвуя в фондировании активов, а каждый рубль выданных кредитов финансируется одним рублём депозитов.

\section{Динамика классов активов}

Для каждого класса $i$ задаются:
\begin{itemize}[noitemsep]
    \item годовой темп амортизации $\alpha_i$ (средний срок кредита $\tau_i = 1/\alpha_i$ лет),
    \item годовой темп новой выдачи $g_{i,t}$ (в долях от текущего объёма портфеля $A_{i,t}$).
\end{itemize}

Шаг модели равен $\Delta t$ (один месяц, $\Delta t = 1/12$).
Динамика объёма портфеля класса $i$ описывается уравнением
\begin{equation}
    A_{i,t+1} = A_{i,t} \bigl(1 - \alpha_i \Delta t + g_{i,t} \Delta t \bigr)
    \equiv A_{i,t} \, m_{i,t},
\end{equation}
где
\[
m_{i,t}=1-\alpha_i\Delta t + g_{i,t}\Delta t
\]
--- валовой коэффициент роста портфеля класса $i$.

Совокупный объём активов равен
\[
A_{t+1} = \sum_i A_{i,t} m_{i,t}.
\]

\section{Ставка новых кредитов и средняя ставка портфеля}

\subsection{Ставка новых выдач}

Новые кредиты класса $i$ выдаются под годовую процентную ставку
\begin{equation}
    R_{i,t} = r_t + s_{i}(g_{i,t}),
\end{equation}
где $r_t$ --- краткосрочная безрисковая ставка (стоимость фондирования), а $s_i(\cdot)$ --- кредитный спред, зависящий от темпа роста $g_{i,t}$ и характеристик соответствующего класса.

Используется линейная спецификация:
\begin{equation}
    s_{i}(g) = s_{0,i} - \kappa_{i} \cdot (g - \alpha_{i}).
\end{equation}
Здесь $s_{0,i}$ --- базовый спред при нормальном обновлении портфеля ($g=\alpha_i$), а $\kappa_i>0$ характеризует чувствительность спреда к объёму новых выдач.

Параметры $s_{0,i}$ и $\kappa_i$ могут зависеть от дюрации $\tau_i$, отражая временную структуру рыночных ставок и премий за ликвидность.

При необходимости на спред могут накладываться нижняя и верхняя границы:
\begin{equation}
    s_{i}(g)=
    \max\bigl(s_{\min},
    \min(s_{\max},
    s_{0,i}-\kappa_i(g-\alpha_i))\bigr).
\end{equation}

\subsection{Динамика средней ставки класса}

Средняя фиксированная ставка портфеля класса $i$, обозначаемая $\overline{R}_{i,t}$, обновляется как взвешенное среднее ставки существующего портфеля и ставки новых выдач:
\begin{equation}
    \overline{R}_{i,t+1}
    =
    \frac{
    (1-\alpha_i\Delta t)\,\overline{R}_{i,t}
    +
    g_{i,t}\Delta t\,R_{i,t}
    }
    {m_{i,t}}.
\end{equation}

\section{Процентный доход и капитал}

Чистый процентный доход за период (в денежных единицах) определяется как
\begin{equation}
    \Pi_t =
    \Bigl(
    \sum_{i=1}^{M}
    A_{i,t}
    (\overline{R}_{i,t}-r_t)
    \Bigr)\Delta t
    +
    r_tK_t\Delta t.
\end{equation}

Капитал до выплаты дивидендов равен
\begin{equation}
    K_t^{\text{pre}}
    =
    K_t+\Pi_t.
\end{equation}

Коэффициент достаточности капитала (Capital Adequacy Ratio, CAR) определяется как
\begin{equation}
    c_t=\frac{K_t}{A_t}.
\end{equation}

\section{Дефолт, дивиденды и терминальная стоимость}

\subsection{Условие дефолта}

Предполагается, что дефолт наступает, если после получения прибыли коэффициент достаточности капитала становится ниже регуляторного минимума $c_0$:
\begin{equation}
    \frac{K_t^{\text{pre}}}{A_{t+1}}<c_0.
\end{equation}

В случае дефолта акционер получает восстановительную стоимость
\begin{equation}
    Rec_t=(1-\delta)\,\pos{K_t^{\text{pre}}},
\end{equation}
где $\delta\in[0,1]$ --- доля потерь при разрешении банка (haircut).

\subsection{Дивидендная политика}

Если дефолт не наступил, дивиденды выплачиваются из положительной прибыли, однако их величина не может снизить капитал ниже целевого уровня $c^*$ относительно нового размера баланса:
\begin{equation}
    Div_t=
    \min\Bigl(
    p_t\,\pos{\Pi_t},
    \;
    \pos{K_t^{\text{pre}}-c^*A_{t+1}}
    \Bigr),
\end{equation}
где $p_t\in[0,1]$ --- управляемый коэффициент выплаты дивидендов (payout ratio).

После выплаты дивидендов капитал становится равным
\begin{equation}
    K_{t+1}=K_t^{\text{pre}}-Div_t.
\end{equation}

\subsection{Терминальное условие}

Если банк продолжает деятельность до момента $T$, терминальная стоимость определяется как
\begin{equation}
    TV_T=K_T+\beta\theta A_T,
\end{equation}
где $\beta\in[0,1]$ --- весовой коэффициент, а $\theta$ --- стоимость единицы активов в стационарном режиме, которая может вычисляться, например, по формуле
\begin{equation}
    \theta=
    \frac{
    \overline{s}_0+\mu c^*
    }{
    \text{cost\_of\_equity}
    },
\end{equation}
где $\overline{s}_0$ --- средний базовый спред кредитного портфеля.

\section{Динамика краткосрочной ставки}

Краткосрочная ставка $r_t$ описывается дискретным процессом Орнштейна--Уленбека (AR(1)):
\begin{equation}
    r_{t+1}
    =
    \mu+\rho(r_t-\mu)+\sigma\varepsilon_{t+1},
    \qquad
    \varepsilon_{t+1}\sim\mathcal{N}(0,1),
\end{equation}
где $\mu$ --- долгосрочное среднее значение ставки, $\rho=e^{-a\Delta t}$, $\sigma$ --- годовая волатильность, а $a$ --- скорость возврата к среднему.

\section{Нормированная задача}

Благодаря линейной однородности модели по масштабу активов задачу можно нормировать на величину $A_t$.

Введём относительный капитал
\[
c_t=\frac{K_t}{A_t},
\]
доли отдельных классов
\[
w_{i,t}=\frac{A_{i,t}}{A_t},
\qquad
\sum_i w_{i,t}=1,
\]
а также средние ставки $\overline{R}_{i,t}$.

Тогда динамика долей имеет вид
\begin{equation}
    w_{i,t+1}
    =
    \frac{w_{i,t}m_{i,t}}
    {\sum_j w_{j,t}m_{j,t}}.
\end{equation}

Прибыль на единицу активов равна
\begin{equation}
    \pi_t
    =
    \frac{\Pi_t}{A_t}
    =
    \Bigl(
    \sum_i
    w_{i,t}
    (\overline{R}_{i,t}-r_t)
    \Bigr)\Delta t
    +
    r_tc_t\Delta t.
\end{equation}

Средний коэффициент роста кредитного портфеля определяется как
\begin{equation}
    \overline{m}_t
    =
    \sum_i
    w_{i,t}m_{i,t}.
\end{equation}

Относительный капитал после выплаты дивидендов:
\begin{equation}
    c_{t+1}
    =
    \frac{c_t+\pi_t-d_t}
    {\overline{m}_t},
\end{equation}
где
\[
d_t=\frac{Div_t}{A_t}
\]
--- нормированный дивиденд,
\begin{equation}
    d_t=
    \min\Bigl(
    p_t\max(\pi_t,0),
    \;
    \max(c_t+\pi_t-c^*\overline{m}_t,0)
    \Bigr).
\end{equation}

Условие дефолта принимает вид
\begin{equation}
    \frac{c_t+\pi_t}{\overline{m}_t}<c_0
    \quad\Rightarrow\quad
    \mathcal{D}_t=1.
\end{equation}

Нормированная восстановительная стоимость равна
\[
(1-\delta)\pos{c_t+\pi_t}.
\]

Терминальное условие (на единицу активов):
\begin{equation}
    v_T(c,\mathbf{w},\overline{\mathbf{R}},r)
    =
    c+\beta\theta.
\end{equation}

\section{Уравнение Беллмана}

Решаем задачу оптимального контроля

\begin{equation}
    \max_{\pi}\; \mathbb{E}\left[
        \sum_{t=0}^{\min(\tau,T)} \gamma^{\,t} d_t
        + \gamma^{\tau} (1-\delta)\max(c_{\tau}+\pi_{\tau},\,0)\,\mathbf{1}_{\{\tau \le T\}}
        + \gamma^{T} \bigl(c_T + \beta\theta\bigr)\,\mathbf{1}_{\{\tau > T\}}
    \right],
\end{equation}

где $\tau$ -- момент дефолта. Запишем для такой задачи уравнение Беллмана.

Пусть
\[
v_t(c,\mathbf{w},\overline{\mathbf{R}},r)
\]
обозначает стоимость банка в расчёте на единицу активов в момент времени $t$.

Оптимальное управление $\{g_{i,t}\}_{i=1}^{M}$ и $p_t$ определяется из уравнения Беллмана:
\begin{equation}
\begin{aligned}
v_t(c,\mathbf{w},\overline{\mathbf{R}},r)
=
\max_{g_1,\dots,g_M,p}
\Bigg\{
&
\1\left\{
\frac{c+\pi}{\overline{m}}<c_0
\right\}
(1-\delta)\max(c+\pi,0)
\\
&
+
\1\left\{
\frac{c+\pi}{\overline{m}}\ge c_0
\right\}
\Bigl(
d
+
\gamma\overline{m}\,
\E\bigl[
v_{t+1}(c',\mathbf{w}',\overline{\mathbf{R}}',r')
\bigr]
\Bigr)
\Bigg\},
\end{aligned}
\end{equation}
где
\[
\gamma=e^{-r_{\text{coe}}\Delta t}
\]
--- одношаговый коэффициент дисконтирования, а величины $\pi$, $\overline{m}$, $d$, $c'$, $w_i'$, $\overline{R}_i'$ и $r'$ вычисляются по формулам, приведённым выше.

\section{Обсуждение размерности и методов решения}

Состояние нормированной задачи включает $M-1$ независимых долей $w_i$, $M$ средних ставок $\overline{R}_i$, а также переменные $c$ и $r$. Таким образом, размерность пространства состояний равна $2M+1$.

При небольшом числе классов ($M\le3$) задачу можно решать методом динамического программирования на сетке.

При больших значениях $M$ требуется использовать методы приближённого динамического программирования (например, нейронные сети или policy gradient), либо вводить дополнительные упрощения, такие как агрегация классов кредитов или фиксация структуры портфеля.

Предложенная модель является естественным обобщением базовой одноклассовой постановки на случай нескольких кредитных классов с различными сроками погашения и позволяет исследовать оптимальное управление дюрацией портфеля, влияние наклона кривой доходности и премий за ликвидность на оптимальную политику банка.

\end{document}